\clearpage
\pdfbookmark[2]{6. Gauss an Olbers, 13. Januar 1815}{parabolisch-006}
\begin{center}
[6.]\par
\medskip
\textsc{Gauss} an \textsc{Olbers}, Göttingen, 13. Januar 1815.
\end{center}
\noindent\hspace*{\parindent}.\ .\ .\ .\ .\ .\par
Ich habe mich diese Woche hindurch weiter mit der Theorie der Kometenbahnen beschäftigt und bin auf eine weitere Vereinfachung gekommen, nach der mir nun nichts mehr zu wünschen übrig zu bleiben scheint. Ich finde nämlich für die Zwischenzeit den Ausdruck
\[
t''-t=3mk\BareRoot{(r+r'')}\cdot\BareRoot{\left(1-\frac{27m^4k^6}{4(t''-t)^4}\right)}\mathpunct{\raisebox{.45ex}{,}}
\]
der in allen in der Praxis vorkommenden Fällen genau genug ist. Ganz streng genommen sollte die letzte Quadratwurzel eine unendliche Reihe sein, die, wenn ich der Kürze halber \(\frac{27m^4k^6}{4(t''-t)^4}=x\) setze, nach meiner Entwicklung \(=\BareRoot{(1-x-2x^3-8x^4-\ldots)}\) wird und in der also glücklicherweise das zweite Glied fehlt. Selbst wenn \(k=\tfrac12(r+r'')\), wäre \(x\) nur etwa \(\frac1{47}\), der Fehler also, \SourceLeaf{342}wenn man sich auf \(\BareRoot{(1-x)}\) beschränkt, nur etwa \(\frac1{100000}\) des Ganzen\HistoricalNoteEight{In einem Exemplar der deutschen Übersetzung der Abhandlung »Observationes Cometae secundi A. 1813 etc.« findet sich die handschriftliche Bemerkung:] NB. Die Reihe \(1-z-2z^3-8z^4-45z^5-272z^6-\ldots\) konvergiert, solange \(z<\frac{27}{256}\). An dieser Grenze, der \(k=r+r''\) entspricht, ist ihr Wert \(=\frac89\).}. Ich wende nun die obige Formel auf folgende Art an. Ich setze
\[
u=A\operatorname{tang}Q,
\]
wodurch
\[
k=\frac A{\cos Q}\quad\text{und}\quad\rho=\frac{g\cos(Q-\varphi)}{h\cos Q}\mathpunct{\cdot}
\]

Man hat also
\[
\begin{aligned}
(t''-t)^2\cos Q^2={}&9mmAA\left(1-\frac{27m^4A^6}{4(t''-t)^4\cos Q^6}\right)\\[4pt]
&\quad\Biggl\{\BareRoot{\left(BB+\left(\frac{A\operatorname{tang}Q+c}{b}\right)^2\right)}\\[4pt]
&\qquad+\BareRoot{\left(B''B''+\frac{A\operatorname{tang}Q+c''}{b''}\right)^2}\Biggr\}
\end{aligned}
\]
oder, wenn ich die folgenden Bezeichnungen einführe,
\[
\begin{gathered}
\frac{9mmAAB}{(t''-t)^2}=\delta,\qquad\frac{9mmAAB''}{(t''-t)^2}=\delta'',\\[5pt]
\frac A{bB}=\epsilon,\qquad\frac A{b''B''}=\epsilon'',\\[5pt]
\frac c{bB}=\eta,\qquad\frac{c''}{b''B''}=\eta'',\qquad\frac{27m^4A^6}{4(t''-t)^4}=\omega,
\end{gathered}
\]
ergibt sich
\[
\begin{aligned}
\frac{\cos Q^2}{1-\dfrac\omega{\cos Q^6}}={}&\delta\BareRoot{\bigl(1+(\epsilon\operatorname{tang}Q+\eta)^2\bigr)}\\[4pt]
&+\delta''\BareRoot{\bigl(1+(\epsilon''\operatorname{tang}Q+\eta'')^2\bigr)}.
\end{aligned}
\]

Etwas Einfacheres ist wohl nicht zu wünschen. Man könnte noch zwei Hilfswinkel \(\vartheta\), \(\vartheta''\) einführen, indem man
\[
\begin{aligned}
\delta\epsilon&=\lambda\sin\vartheta, &\delta''\epsilon''&=\lambda''\sin\vartheta''\\
\delta\eta&=\lambda\cos\vartheta, &\delta''\eta''&=\lambda''\cos\vartheta''
\end{aligned}
\]
setzte, wodurch die Gleichung folgende Gestalt erhielte
\[
\begin{aligned}
\frac{\cos Q^3}{1-\dfrac\omega{\cos Q^6}}={}&\BareRoot{\bigl(\delta\delta\cos Q^2+\lambda\lambda\cos(Q-\vartheta)^2\bigr)}\\[4pt]
&+\BareRoot{\bigl(\delta''\delta''\cos Q^2+\lambda''\lambda''\cos(Q-\vartheta'')^2\bigr)}.
\end{aligned}
\]

Man könnte sogar noch weiter gehen und der Gleichung durch bekannte Umformungen die folgende Gestalt geben.
\[
\begin{aligned}
\frac{\cos Q^3}{1-\dfrac\omega{\cos Q^6}}={}&\mu\BareRoot{\bigl(1+\nu\nu\cos(Q+\pi)^2\bigr)}\\[4pt]
&+\mu''\BareRoot{\bigl(1+\nu''\nu''\cos(Q+\pi'')^2\bigr)}.
\end{aligned}
\]

\SourceLeaf{343}Doch der Gewinn\footnote{Er würde nämlich \textls[150]{nur} darin bestehen, dass man sofort das Maximum und das Minimum der Werte der beiden Wurzelausdrücke beurteilen könnte. Vernachlässigte man zunächst \(\omega\), so erhielte man dadurch \textls[150]{Grenzen} für \(\cos Q^3\) und folglich für \(Q\). Bei der Berechnung jedes einzelnen Versuchs muss man immer noch 7-mal nachschlagen, gleich welche Form man wählt.} wäre zu gering, als dass sich die Mühe dieser Umformungen lohnte. Die Gleichung ist stets vollkommen exakt, sobald man an die Stelle des Nenners auf der linken Seite die Reihe setzt
\[
1-\frac\omega{\cos Q^6}-\frac{2\omega^3}{\cos Q^{18}}-\frac{8\omega^4}{\cos Q^{24}}-\text{usw.}
\]

Das Problem ist also auf eine Gleichung mit sieben bekannten Größen zurückgeführt, und ich glaube nicht, dass es sich auf eine geringere Anzahl reduzieren lässt.

Die allgemeine Aufgabe stellt sich, wenn der mittlere Ort unvollständig bekannt ist und man nur irgendeinen Großkreis kennt, auf dem er liegt. Die vorliegende Aufgabe ist eigentlich nur der einzelne Fall, in dem dieser Großkreis durch den mittleren \(\odot\)ort geht. Man muss auf die allgemeine Aufgabe zurückgreifen, sooft dieser \(\odot\)ort nahe bei der Richtung der geoz[entrischen] Bewegung liegt. Diese allgemeine Aufgabe habe ich nun auch auf analoge Art durchgeführt und ein Musterbeispiel sorgfältig berechnet\HistoricalNoteEight{Siehe die folgende Notiz.]}. Der Natur der Sache nach ist die Arbeit hier viel verwickelter. Ich komme am Ende auf zwei Gleichungen mit 11\HistoricalNoteDoubleUnbracketed{Für jeden Versuch muss man 11-mal nachschlagen.} bekannten Größen und zweifle auch, ob sich diese Anzahl verringern lässt. Wenn die Berechnung einer Kometenbahn im ersten Fall nach allen Abkürzungen jetzt mindestens 1 Stunde erfordert (so ist es etwa bei mir), könnte der andere Fall leicht 2—3 Stunden erfordern.
\par\medskip\noindent\makebox[\linewidth]{\rule{.18\linewidth}{.35pt}}\par
